% ============================================================
% CUMCM 国赛 C 题评价类模板 — v7.2
% 编译:xelatex main.tex(跑两遍)
% 金标准:六段子结构 / 三段式公式 / 四重检验 / 算法对比 / 创新量化
% 核心强化:指标体系 / 权重计算 / 一致性检验
% ============================================================
\documentclass[12pt,a4paper]{ctexart}

\usepackage{amsmath,amssymb,amsfonts}
\usepackage{graphicx}
\usepackage{booktabs}
\usepackage{geometry}
\usepackage{float}
\usepackage{caption}
\usepackage{subcaption}
\usepackage{listings}
\usepackage{hyperref}
\usepackage{enumitem}
\usepackage{xcolor}
\usepackage{indentfirst}
\usepackage{titlesec}
\usepackage{fancyhdr}
\usepackage{setspace}
\usepackage{tikz}

\geometry{top=2.5cm, bottom=2.5cm, left=2.5cm, right=2.5cm, headheight=14pt}
% ============ 排版规范（对齐国赛优秀论文 B159） ============
% 行距 1.5 倍，首行缩进 2 字符
\linespread{1.5}\selectfont
\setlength{\parindent}{2em}

% 章节标题居中，二级标题左对齐
\titleformat{\section}{\centering\large\bfseries\heiti}{\thesection}{0.8em}{}
\titleformat{\subsection}{\normalsize\bfseries\heiti}{\thesubsection}{0.8em}{}
\titleformat{\subsubsection}{\normalsize\bfseries\heiti}{\thesubsubsection}{0.6em}{}

% 页眉空、页脚页码居中
\pagestyle{fancy}
\fancyhf{}
\fancyfoot[C]{\small\thepage}
\renewcommand{\headrulewidth}{0pt}


\setCJKmainfont{SimSun}[BoldFont=SimHei, ItalicFont=KaiTi]
\setCJKsansfont{SimHei}
\setCJKmonofont{FangSong}

% 图/表题自动编号前缀（ctexart 默认不设置，缺失会导致图题无「图」字）
\renewcommand{\figurename}{图}
\renewcommand{\tablename}{表}


\DeclareMathOperator{\erf}{erf}
\DeclareMathOperator{\sgn}{sgn}

\hypersetup{colorlinks=true, linkcolor=black, citecolor=black, urlcolor=blue}

\lstset{
  basicstyle=\ttfamily\small,
  breaklines=true,
  frame=single,
  numbers=left,
  numberstyle=\tiny,
  keywordstyle=\color{blue},
  commentstyle=\color{green!50!black},
}

\title{\textbf{[评价对象]的综合评价体系建模与求解}}
\author{队伍编号:\textbf{XXXX}}
\date{}

\begin{document}

\maketitle

% ============================================================
\begin{abstract}
\noindent
\textbf{第一段}:行业背景 + 评价对象 + 三层框架(指标-权重-排序)。

\textbf{第二段}:问题一构建 [N 项] 指标体系,采用 AHP+熵权组合赋权,TOPSIS 排序得到最优方案 [方案名]($C_i = 0.87$);问题二...;问题三...

\textbf{第三段}:3 项创新(指标体系创新 / 组合赋权 / 一致性提升 $X\%$)+ 四重检验指标 + 跨行业迁移。

\vspace{0.5em}
\noindent\textbf{关键词}:[对象] [评价] [AHP] [TOPSIS] [一致性]
\end{abstract}

\newpage
\newpage

% ============================================================
\section{问题重述}
\subsection{问题背景}
[背景:聚焦评价对象与决策需求]
\subsection{问题提出}
\begin{itemize}
  \item 问题一:[指标体系构建]
  \item 问题二:[权重计算与排序]
  \item 问题三:[方案择优与灵敏度]
\end{itemize}

% ============================================================
\section{问题分析}
\begin{figure}[H]
\centering
\includegraphics[width=0.85\textwidth]{../figures/pdf/fig1_indicator_system.pdf}
\caption{图1 评价指标体系(目标层-准则层-指标层三层结构)}
\label{fig:indicator}
\end{figure}

% ============================================================
\section{模型假设}
\begin{enumerate}
  \item 假设1:评价指标相互独立(或相关性已处理)
  \item 假设2:专家打分具有一致性
  \item 假设3:数据来源真实可靠
  \item 假设4:无量纲化不改变指标相对优劣
\end{enumerate}

% ============================================================
\section{符号说明}
\begin{table}[H]
\centering
\caption{符号说明}
\label{tab:symbols}
\begin{tabular}{llll}
\toprule
符号 & 含义 & 取值范围 & 单位 \\
\midrule
$X_{ij}$ & 第 $i$ 方案第 $j$ 指标原始值 & - & - \\
$x_{ij}$ & 无量纲化后值 & $[0, 1]$ & - \\
$w_j$ & 第 $j$ 指标权重 & $[0, 1]$ & - \\
$C_i$ & 第 $i$ 方案贴近度 & $[0, 1]$ & - \\
\bottomrule
\end{tabular}
\end{table}

% ============================================================
\section{模型建立与求解}

\subsection{问题一:指标体系构建}

\subsubsection{参数配置与数据预处理}
[数据来源、字段说明、异常处理]

\subsubsection{基础经典模型}
[基础指标体系,如等权重处理]
\begin{equation}
C_i = \frac{1}{n} \sum_{j=1}^{n} x_{ij}
\label{eq:equal_weight}
\end{equation}
其中 $x_{ij}$ 为无量纲化值,$n$ 为指标数。式~\eqref{eq:equal_weight} 假设各指标等权重,未考虑指标重要性差异。

\subsubsection{传统模型缺陷剖析}
\begin{enumerate}
  \item 缺陷1:等权重忽略指标重要性差异
  \item 缺陷2:未区分正向/负向指标,排序失真 $X\%$
  \item 缺陷3:主观赋权缺乏客观性
\end{enumerate}

\subsubsection{创新改进模型:组合赋权}
本文构建 AHP+熵权组合赋权模型:
\paragraph{AHP 主观权重}
\begin{equation}
\mathbf{A} \mathbf{w}_{\text{AHP}} = \lambda_{\max} \mathbf{w}_{\text{AHP}}
\label{eq:ahp}
\end{equation}
其中 $\mathbf{A}$ 为判断矩阵,$\lambda_{\max}$ 为最大特征值。一致性比率 $\text{CR} = \text{CI}/\text{RI} < 0.1$。

\paragraph{熵权客观权重}
\begin{equation}
w_{\text{entropy},j} = \frac{1 - E_j}{\sum_{k=1}^{n} (1 - E_k)}, \quad E_j = -\frac{1}{\ln m} \sum_{i=1}^{m} p_{ij} \ln p_{ij}
\label{eq:entropy}
\end{equation}

式~\eqref{eq:ahp} 给出主观权重,式~\eqref{eq:entropy} 给出客观权重。

\paragraph{组合权重}
\begin{equation}
w_j = \alpha w_{\text{AHP},j} + (1 - \alpha) w_{\text{entropy},j}, \quad \alpha = 0.5
\label{eq:combined}
\end{equation}

式~\eqref{eq:combined} 给出组合权重,其中 $\alpha = 0.5$ 平衡主客观信息。

\subsubsection{算法设计与求解:TOPSIS 排序}

\paragraph{三算法对比表}
\begin{table}[H]
\centering
\caption{权重计算方法对比}
\label{tab:algo_compare}
\begin{tabular}{lccccc}
\toprule
方法 & 客观性 & 一致性 CR & 区分度 & 计算时间 (s) & 评级 \\
\midrule
AHP+熵权+TOPSIS & 高 & 0.087 & 0.42 & 0.8 & ★★★ \\
纯 AHP & 低 & 0.092 & 0.31 & 0.5 & ★★☆ \\
纯熵权 & 高 & - & 0.35 & 0.3 & ★★☆ \\
\bottomrule
\end{tabular}
\end{table}

选用 AHP+熵权组合的理由:客观性最高,一致性 CR=0.087 < 0.1,区分度最大($0.42$)。

\subsubsection{数值结果与可视化}

\begin{figure}[H]
\centering
\includegraphics[width=0.85\textwidth]{../figures/pdf/fig2_radar.pdf}
\caption{图2 各方案多维度雷达对比(方案 A 各维度均衡,$C_i = 0.87$ 最优)}
\label{fig:radar}
\end{figure}

\begin{table}[H]
\centering
\caption{问题一综合评价结果}
\label{tab:result1}
\begin{tabular}{lcccc}
\toprule
方案 & 组合权重得分 & TOPSIS 贴近度 $C_i$ & 排名 & 较传统方法 \\
\midrule
方案 A & 0.87 & 0.87 & 1 & +0.12 \\
方案 B & 0.78 & 0.78 & 2 & +0.08 \\
方案 C & 0.65 & 0.65 & 3 & +0.05 \\
\bottomrule
\end{tabular}
\end{table}

图~\ref{fig:radar} 显示方案 A 在精度、效率、稳定性维度均领先,综合贴近度 $C_i = 0.87$ 排名第一,较传统等权重方法提升 $0.12$。

\subsection{问题二}
% 同六段子结构

\subsection{问题三}
% 同六段子结构

% ============================================================
\section{模型检验}

\subsection{拟合精度检验}
评价指标的区分度 $\text{CV} = 0.18$,Kendall 一致性系数 $\tau = 0.92$($p < 0.01$)。

\subsection{灵敏度分析}
\begin{table}[H]
\centering
\caption{权重灵敏度分析}
\label{tab:sensitivity}
\begin{tabular}{lccccc}
\toprule
权重 & 基准值 & ±10\% 影响 & ±20\% 影响 & 弹性系数 $\varepsilon$ & 分级 \\
\midrule
$w_1$ & 0.25 & ±0.8\% & ±1.6\% & 0.08 & 低 \\
$w_2$ & 0.30 & ±2.5\% & ±5.2\% & 0.26 & 中 \\
$w_3$ & 0.45 & ±4.1\% & ±8.7\% & 0.41 & 中 \\
\bottomrule
\end{tabular}
\end{table}

\begin{figure}[H]
\centering
\includegraphics[width=0.85\textwidth]{../figures/pdf/figN_weight_sensitivity.pdf}
\caption{图N 权重灵敏度热力图($w_3$ 为中灵敏度参数)}
\label{fig:weight_sens}
\end{figure}

\subsection{蒙特卡洛验证}
$N_{\text{sim}} = 200$ 次随机扰动权重,排序一致性 $92.5\%$,$95\%\text{CI} = [0.85, 0.91]$。

\subsection{假设误差量化}
\begin{table}[H]
\centering
\caption{假设误差量化}
\label{tab:assumption_error}
\begin{tabular}{lcc}
\toprule
假设 & 偏差量化 & 对排序影响 \\
\midrule
假设1:指标独立 & 相关性 $\rho = 0.32$ & 排序变动率 $5\%$ \\
假设2:专家一致 & CR=0.087 & 排序稳定 \\
假设3:数据真实 & 异常值 $2.1\%$ & 排序变动率 $3\%$ \\
\bottomrule
\end{tabular}
\end{table}

% ============================================================
\section{模型优缺点与改进}
\textbf{创新点}:
\begin{enumerate}
  \item 构建 [N 项] 指标体系,覆盖 [维度] 较传统方法完整性提升 $35\%$
  \item AHP+熵权组合赋权,客观性较纯 AHP 提升 $28\%$
  \item TOPSIS 排序一致性 $\tau = 0.92$,较单一方法提升 $15\%$
\end{enumerate}

% ============================================================
\section{参考文献}
\begin{thebibliography}{99}
\bibitem{ref1} 作者. 标题[J]. 期刊, 年份, 卷(期): 页码.
\bibitem{ref2} Author A. Title[J]. Journal, Year, Vol: Pages.
% 近5年 ≥40%, 外文 ≥30%, 总数 ≥10
\end{thebibliography}

% ============================================================
\appendix
\section{附录1 数据说明}
[指标原始数据样表]

\section{附录2 代码清单}
% ============ 国一铁律：附录必须含完整代码，缺代码取消评奖资格 ============
% 完整代码用 \lstinputlisting 从 code/ 目录读入，禁止只贴片段/伪代码。
% 每个脚本附输入/输出/功能说明，确保可复现（全部固定 SEED=42）。
\begin{table}[H]
\centering
\caption{代码模块说明（完整代码见附件，固定 SEED=42 可复现）}
\label{tab:code_modules}
\begin{tabular}{llll}
\toprule
脚本 & 功能 & 输入 & 输出 \\
\midrule
\texttt{code/00\_data\_prep.py} & 数据预处理与质量校验 & 附件1.xlsx & \texttt{data/clean/*.csv} \\
\texttt{code/01\_q1\_analysis.py} & 问题一求解 & 清洗后数据 & \texttt{results/q1\_*.json} \\
\texttt{code/02\_q2\_classify.py} & 问题二求解 & 关键词表 & \texttt{results/q2\_*.json} \\
\texttt{code/03\_q3\_optimize.py} & 问题三求解 & 问题二结果 & \texttt{results/q3\_*.json} \\
\texttt{code/04\_q4\_robust.py} & 问题四求解 & 问题三模型 & \texttt{results/q4\_*.json} \\
\texttt{code/05\_figures.py} & 图表生成 & \texttt{results/*.json} & \texttt{figures/png}, \texttt{figures/pdf} \\
\bottomrule
\end{tabular}
\end{table}

% 逐文件引入完整代码（示例）：
% \lstinputlisting[language=Python, caption=code/01_q1_analysis.py]{../code/01_q1_analysis.py}


\section{附录3 公式推导}
[组合赋权完整推导]

\end{document}
